\section{Quadratic moments and certificate systems}\label{sec:setting}

\subsection{The full hull and the positive-loop hull}\label{set:cones}
Let $C_n=[0,1]^n$. We write a quadratic as
\begin{equation}\label{eq:set:quadratic}
 q(x)=q_0+\sum_{i=1}^n a_ix_i+\sum_{i=1}^n b_ix_i^2
                  +\sum_{1\le i<j\le n}c_{ij}x_ix_j.
\end{equation}
Thus $c_{ij}$ is the coefficient of the monomial $x_ix_j$, rather than an
off-diagonal entry of a symmetric matrix representing $q$. Define
\[
 \mathcal P_n=\{q:q(x)\ge0\text{ for every }x\in C_n\},\qquad
 \mathcal P_n^+=\{q\in\mathcal P_n:b_i\ge0\text{ for every }i\}.
\]
The corresponding moment sets are
\begin{equation}\label{eq:set:hulls}
 \mathcal Q_n=\conv\{(x,xx^\top):x\in C_n\},\qquad
 \mathcal H_n=\mathcal Q_n+
       \{(0,\diag(s)):s\in\R_+^n\}.
\end{equation}
The set $\mathcal Q_n$ is compact. The set $\mathcal H_n$ is closed because
it is the sum of a compact set and a closed cone. A point $(m,Y)$ in
$\mathcal Q_n$ records the first and second moments of a finite probability
law on the box. In $\mathcal H_n$, each diagonal second moment may also
contain nonnegative slack. We use the name \emph{positive-loop hull} for this
second set: a loop is the graph coordinate representing a square, and its
sign specifies the permitted direction of diagonal slack.

For a quadratic, its evaluation at moments is
\begin{equation}\label{eq:set:pairing}
 \mathcal L_{m,Y}(q)=q_0+\sum_i a_im_i+\sum_i b_iY_{ii}
                         +\sum_{i<j}c_{ij}Y_{ij}.
\end{equation}
Every $q\in\mathcal P_n$ gives a valid inequality $\mathcal L_{m,Y}(q)\ge0$ for
$\mathcal Q_n$. It is valid for $\mathcal H_n$ exactly when $q\in\mathcal P_n^+$.
Indeed, evaluating at $(x,xx^\top)$ gives $q(x)$, while adding diagonal
slack changes the value by $\sum_i b_is_i$. Conversely, finite-dimensional
separation shows that these valid inequalities describe the two closed
convex sets.

These are joint hulls of all recorded products. A lift of an epigraph for
one fixed objective need not describe the full moment hull. This distinction
will matter in the nonrepresentability results.

\subsection{Normalization and duality}\label{set:dual}
Taking cone duals requires retaining the constant coordinate. Let
\[
 \widehat{\mathcal Q}_n
   =\cone\{(1,x,xx^\top):x\in C_n\},\qquad
 \widehat{\mathcal H}_n
   =\widehat{\mathcal Q}_n+
       \{(0,0,\diag(s)):s\in\R_+^n\}.
\]
Both cones are closed. In the first cone, every coordinate is bounded in
absolute value by its mass coordinate, so a convergent sequence with mass
tending to zero has limit zero. The same argument in the second cone leaves
only the displayed diagonal slack when the mass tends to zero. For positive
mass, closure follows by rescaling to the compact set $\mathcal Q_n$ and
extracting a convergent slack vector.

Under the pairing obtained from \eqref{eq:set:pairing} by replacing $q_0$
with $tq_0$, their dual cones are $\mathcal P_n$ and $\mathcal P_n^+$,
respectively. Their mass-one slices are $\mathcal Q_n$ and $\mathcal H_n$.
The zero-mass slack in $\widehat{\mathcal H}_n$ cannot be discarded when
using homogeneous duality.

We call a set a \emph{spectrahedral shadow}, or equivalently a set with a
\emph{finite semidefinite lift}, if it can be written as
\begin{equation}\label{eq:set:shadow}
 \{y:\exists z,\ A_0+\textstyle\sum_i y_iA_i+\sum_j z_jB_j\succeq0\},
\end{equation}
possibly with additional affine equalities. Several PSD blocks are allowed,
since they can be combined in one block-diagonal matrix. The word ``finite''
places no polynomial bound on the number of auxiliaries or the block orders.
Consequently, failure of finite representability is stronger than failure
of a particular small relaxation. It is not a computational hardness theorem.

\subsection{The three-variable disjoint-support system}\label{set:disjoint}
For the remainder of the local certificate analysis, $n=3$. Let $W$ be
the ten-dimensional vector space of quadratics, including their constants.
Let $V$ be the twenty-dimensional span of the monomials $x^\alpha$ with
$\alpha\in\{0,1,2\}^3$ and at most one exponent equal to two. For disjoint
subsets $A,B\subseteq\{1,2,3\}$, put
\[
 w_{A,B}(x)=\prod_{i\in A}x_i\prod_{j\in B}(1-x_j),\qquad
 F_{A,B}=\{1,2,3\}\setminus(A\cup B).
\]
Let $v_F(x)$ list the constant one and the coordinates indexed by $F$.
Define the cone
\begin{equation}\label{eq:set:disjoint}
 \mathcal K_D=\left\{\sum_{A\cap B=\varnothing}
       w_{A,B}(x)v_{F_{A,B}}(x)^\top G_{A,B}v_{F_{A,B}}(x):
       G_{A,B}\succeq0\right\}\subseteq V,
 \qquad \mathcal D_3=\mathcal K_D\cap W.
\end{equation}
Each summand is nonnegative on the cube. The matrix representation permits
any finite sum of affine squares on the unused coordinates. The requirement
that the final sum belongs to $W$ imposes cancellation of all monomials of
degree greater than two; it must not be replaced by a degree restriction on
each summand. In particular, $\mathcal D_3\subseteq\mathcal P_3^+$: the coefficient
of $x_i^2$ comes only from a nonnegative diagonal entry of a block whose
unused coordinates include $i$.

The dual local relaxation consists of restrictions to $W$ of functionals
$\ell\in V^*$ satisfying
\begin{equation}\label{eq:set:localizing}
 \ell(1)=1,\qquad
 M_{A,B}(\ell):=
       \ell\bigl(w_{A,B}v_{F_{A,B}}v_{F_{A,B}}^\top\bigr)\succeq0
       \quad(A\cap B=\varnothing).
\end{equation}
Here $\ell$ acts entrywise on a polynomial matrix. Write $\mathcal R_D$ for the
resulting set of quadratic moment pairs $(m,Y)$, with
$m_i=\ell(x_i)$ and $Y_{ij}=\ell(x_ix_j)$. There are $27$ blocks:
eight of order one, twelve of order two, six of order three, and one of
order four. The higher moments in $V\setminus W$ are auxiliary coordinates.
They need not come from one probability law. Every point of $\mathcal H_3$
has such an extension; conversely, equality of the relaxation with that
hull is the substantive exactness question. \Cref{sec:three} proves the
closedness and extension statements needed to use this primal--dual
description without a hidden closure assumption.

\subsection{Diagonal caps and cube symmetries}\label{set:caps}
The full box hull satisfies the three inequalities
\begin{equation}\label{eq:set:caps}
 Y_{ii}\le m_i\quad(i=1,2,3),
\end{equation}
which are the moment forms of $x_i(1-x_i)\ge0$. We call them
\emph{diagonal caps}. They have negative square coefficients and are not
valid for the unbounded positive-loop hull. Thus a proposed description of
$\mathcal Q_3$ and a proposed description of $\mathcal H_3$ are different
statements, even when all other inequalities are the same.

\label{set:symmetry}
Coordinate permutations and complements $x_i\mapsto1-x_i$ preserve the
cube. They induce affine maps on normalized first and second moments and
linear maps on homogeneous moments. Complementing coordinate $i$, for
example, sends
\[
 m_i\mapsto1-m_i,\quad
 Y_{ii}\mapsto1-2m_i+Y_{ii},\quad
 Y_{ij}\mapsto m_j-Y_{ij}\quad(j\ne i),
\]
with simultaneous substitution when more than one coordinate is
complemented. These symmetries preserve $\mathcal Q_3$, $\mathcal H_3$,
and the disjoint-support certificate and moment systems. They provide the
different orientations of the inequality family constructed below.
